\chapter{The Singlet, Bell's Inequality, and No-Cloning}

\lectureoverview{A relative minus sign turns a merely correlated two-spin state into the rotationally invariant singlet. We identify that state by calculating its total spin, use its perfect anti-correlation to formulate a Bell experiment, and translate classical propositions into quantum projection operators. The resulting probabilities violate Bell's classical counting inequality. Questions about locality, experimental timing, and uncertainty then lead to a second consequence of linear quantum mechanics: an unknown state cannot be copied by a universal cloning machine.}

\section{What Difference Can a Sign Make?}

Consider two electrons and suppress the tensor-product symbol in the basis states:
\begin{equation}
  \ket{ud}\equiv\ket{u}_1\otimes\ket{d}_2,
  \qquad
  \ket{du}\equiv\ket{d}_1\otimes\ket{u}_2.
\end{equation}
Two especially simple superpositions differ only by one sign,
\begin{equation}
  \ket{\psi_\pm}
  =
  \frac{\ket{ud}\pm\ket{du}}{\sqrt2}.
  \label{eq:l05-two-signs}
\end{equation}
The factor \(1/\sqrt2\) makes the total probability one. In either state, a measurement along the third axis finds opposite answers: if electron 1 is up, electron 2 is down, and conversely. The sign therefore cannot be detected by that correlation alone. Nevertheless, the two states are physically different.

The calculation will repeatedly use the Pauli actions
\begin{align}
  \sigma_1\ket{u}&=\ket{d},&
  \sigma_1\ket{d}&=\ket{u},\nonumber\\
  \sigma_2\ket{u}&=i\ket{d},&
  \sigma_2\ket{d}&=-i\ket{u},\nonumber\\
  \sigma_3\ket{u}&=\ket{u},&
  \sigma_3\ket{d}&=-\ket{d}.
  \label{eq:l05-pauli-actions}
\end{align}
We reserve \(\sigma_i\) for operators acting on electron 1 and introduce an identical family \(\tau_i\) acting on electron 2.

\begin{figure}[htbp]
\centering
\includegraphics[width=0.62\linewidth]{lecture_05_figure_02.png}
\caption{The complete Pauli-action table used throughout the comparison.}
\lectureframe{00:03:15}
\end{figure}

Neither constituent spin has a preferred direction in either state:
\begin{equation}
  \expect{\sigma_i}_{\psi_\pm}=0,
  \qquad
  \expect{\tau_i}_{\psi_\pm}=0,
  \qquad i=1,2,3.
  \label{eq:l05-local-zero}
\end{equation}
For example, the explicit check requested for the plus state gives
\begin{align}
  \bra{\psi_+}\sigma_3\ket{\psi_+}
  &=
  \frac12
  \bigl(\bra{ud}+\bra{du}\bigr)
  \bigl(\ket{ud}-\ket{du}\bigr)\nonumber\\
  &=
  \frac12(1-1)=0.
  \label{eq:l05-sigma3-expectation}
\end{align}
The cross terms vanish because \(\braket{ud}{du}=0\).

\begin{figure}[htbp]
\centering
\includegraphics[width=0.78\linewidth]{lecture_05_figure_03.png}
\caption{The plus-state expectation value is evaluated by acting with \(\sigma_3\) and using orthogonality.}
\lectureframe{00:07:00}
\end{figure}

\begin{classroomqa}
\lecturetimestamp{00:07:13}
\audiencequestion{Does a zero expectation value mean that a single spin can be measured to have value zero?}
\lecturerresponse{No. A component of one spin has only the outcomes \(+1\) and \(-1\); \(\sigma_i\) has no zero eigenvalue. Equation~\eqref{eq:l05-local-zero} says that repeated measurements are a fair coin toss, not that any individual result is zero.}
\end{classroomqa}

The first plausible diagnostic has therefore failed. Local averages cannot see the relative sign, so we must measure a property of the pair.

\section{The Total Spin Selects the Singlet}

The pair's magnetic moment is represented by the vector sum
\begin{equation}
  \boldsymbol{\Sigma}
  =
  \boldsymbol{\sigma}+\boldsymbol{\tau}.
\end{equation}
Its third component annihilates both basis vectors:
\begin{align}
  (\sigma_3+\tau_3)\ket{ud}
  &=(1-1)\ket{ud}=0,\nonumber\\
  (\sigma_3+\tau_3)\ket{du}
  &=(-1+1)\ket{du}=0.
\end{align}
Consequently,
\begin{equation}
  (\sigma_3+\tau_3)\ket{\psi_\pm}=0.
\end{equation}
Both sign choices have total third component zero on every trial.

The first component separates them. From Equation~\eqref{eq:l05-pauli-actions},
\begin{align}
  (\sigma_1+\tau_1)\ket{ud}
  &=\ket{dd}+\ket{uu},\nonumber\\
  (\sigma_1+\tau_1)\ket{du}
  &=\ket{uu}+\ket{dd}.
\end{align}
Thus
\begin{align}
  (\sigma_1+\tau_1)\ket{\psi_+}
  &=
  \sqrt2\bigl(\ket{uu}+\ket{dd}\bigr),\nonumber\\
  (\sigma_1+\tau_1)\ket{\psi_-}
  &=0.
  \label{eq:l05-total-x}
\end{align}
The plus state may have zero expectation value for this observable, but it is not its zero-eigenvalue eigenstate.

\begin{classroomqa}
\lecturetimestamp{00:13:02}
\audiencequestion{What outcomes can the total component have when the state is not annihilated?}
\lecturerresponse{Each constituent answer is \(+1\) or \(-1\), so their sum can be \(2\), \(0\), or \(-2\). Zero expectation does not force zero on every run.}
\end{classroomqa}

The second component completes the test. Its factors of \(i\) cancel in the antisymmetric combination:
\begin{equation}
  (\sigma_2+\tau_2)\ket{\psi_-}=0.
\end{equation}
All three components therefore annihilate the minus state,
\begin{equation}
  \bigl(\boldsymbol{\sigma}+\boldsymbol{\tau}\bigr)
  \cdot\hat n\,\ket{\psi_-}=0
  \qquad\text{for every unit vector }\hat n.
  \label{eq:l05-singlet-any-axis}
\end{equation}

\begin{figure}[htbp]
\centering
\includegraphics[width=0.84\linewidth]{lecture_05_figure_06.png}
\caption{The component-by-component cancellations that identify the minus state as total spin zero.}
\lectureframe{00:16:30}
\end{figure}

The antisymmetric state \(\ket{\psi_-}\) is the \emph{singlet}. It is rotationally invariant: every axis gives the same total-spin answer, zero. The symmetric state in Equation~\eqref{eq:l05-two-signs} belongs to the triplet sector.

\begin{classroomqa}
\lecturetimestamp{00:19:06}
\audiencequestion{Is the opposite-spin result true only on average?}
\lecturerresponse{For the singlet it is true on every trial. Equation~\eqref{eq:l05-singlet-any-axis} is an eigenvector equation, so measuring the same component on both particles always gives answers whose sum is the eigenvalue zero.}
\end{classroomqa}

\begin{classroomqa}
\lecturetimestamp{00:19:36}
\audiencequestion{Why is the other state called a triplet?}
\lecturerresponse{That name comes from the three states in the total-spin-one multiplet. Its derivation is deferred; the present calculation needs only the contrast between the symmetric state and the one-dimensional singlet sector.}
\end{classroomqa}

If Bob measures any component of his electron, he immediately knows what Alice would obtain by measuring the same component of hers: her answer must be the opposite. This is an Einstein--Podolsky--Rosen correlation. A pair of hidden coins can also carry opposite results, but only for one prearranged question. The singlet supplies perfect anti-correlation for whichever spatial axis is chosen.

\begin{classroomqa}
\lecturetimestamp{00:20:55}
\audiencequestion{Did Einstein, Podolsky, and Rosen first discover entanglement?}
\lecturerresponse{No. Schrödinger discussed the phenomenon and introduced the language of entanglement. The EPR argument made the distant-correlation problem especially sharp.}
\end{classroomqa}

\begin{classroomqa}
\lecturetimestamp{00:23:50}
\audiencequestion{Does the immediate correlation violate causality, and does it matter which observer measures first?}
\lecturerresponse{No usable message is transmitted. Alice's local outcomes remain fifty-fifty whether or not Bob measures, in every inertial ordering. She can verify the correlation only after receiving Bob's result through an ordinary classical channel, which cannot outrun light.}
\end{classroomqa}

\section{Bell's Inequality Is Classical Counting}

The status of Bell's result can feel unstable. On some days it appears to be the deepest feature of quantum mechanics; on others it appears to be one more inevitable consequence of the formalism. Bell himself, and Feynman as well, could express both attitudes. The inequality, however, is elementary. The surprise lies in its violation.

Take an ordinary set of objects. Each object may possess or fail to possess three properties \(A\), \(B\), and \(C\). Let \(N(A,\neg B)\) denote the number that have \(A\) but not \(B\). Classical set logic implies
\begin{equation}
  N(A,\neg B)+N(B,\neg C)
  \geq
  N(A,\neg C).
  \label{eq:l05-bell-count}
\end{equation}
The proof is the set inclusion
\begin{equation}
  A\cap C^{c}
  \subseteq
  (A\cap B^{c})\cup(B\cap C^{c}).
  \label{eq:l05-set-inclusion}
\end{equation}
Indeed, an element of \(A\cap C^c\) either belongs to \(B\) or it does not. In the first case it belongs to \(B\cap C^c\); in the second it belongs to \(A\cap B^c\). Counting the union can only overcount the set on the left.

\begin{figure}[htbp]
\centering
\includegraphics[width=0.78\linewidth]{lecture_05_figure_07.png}
\caption{The board proof partitions three overlapping classical subsets into nonnegative regions.}
\lectureframe{00:31:15}
\end{figure}

Equivalently, label the disjoint regions of the three-set diagram by nonnegative integers \(N_i\). Both sides of Equation~\eqref{eq:l05-bell-count} contain the regions required by \(A\cap C^c\), while the left side contains additional nonnegative regions. Nothing quantum has entered. This is a theorem about subsets and ordinary logical alternatives.

\section{Turning the Set Statement into a Spin Experiment}

Prepare many independent electron pairs in the singlet state and define three propositions about electron 1:
\begin{align}
  A&:\ \text{spin 1 is up at }0^\circ\text{, along the third axis},\nonumber\\
  B&:\ \text{spin 1 is up at }45^\circ\text{ in the }zx\text{-plane},\nonumber\\
  C&:\ \text{spin 1 is up at }90^\circ\text{, along the first axis}.
  \label{eq:l05-abc}
\end{align}

\begin{figure}[htbp]
\centering
\includegraphics[width=0.82\linewidth]{lecture_05_figure_08.png}
\caption{The three spin propositions are fixed at \(0^\circ\), \(45^\circ\), and \(90^\circ\) in one plane.}
\lectureframe{00:39:00}
\end{figure}

In the singlet, the negation of a proposition about electron 1 may be rewritten as a positive proposition about electron 2. For example,
\begin{align}
  \neg B
  &:\ \text{spin 1 is down at }45^\circ\nonumber\\
  &\Longleftrightarrow
  \text{spin 2 is up at }45^\circ.
\end{align}
Likewise, \(\neg C\) becomes the statement that electron 2 is up at \(90^\circ\).

\begin{classroomqa}
\lecturetimestamp{00:40:04}
\audiencequestion{Should \(\neg C\) make electron 2 down at \(90^\circ\)?}
\lecturerresponse{No. First negate \(C\): electron 1 is down at \(90^\circ\). Perfect singlet anti-correlation then makes electron 2 up along that same axis. The two reversals must be kept separate.}
\end{classroomqa}

For a large ensemble, counts and probabilities differ only by the common number of prepared pairs. Rotational invariance also gives
\begin{equation}
  P(A,\neg B)=P(B,\neg C),
\end{equation}
because the second experiment is the first rotated by \(45^\circ\). Bell's classical inequality would therefore demand
\begin{equation}
  2P(A,\neg B)\geq P(A,\neg C).
  \label{eq:l05-bell-probability}
\end{equation}
We now need an efficient way to calculate these probabilities.

\section{Propositions Become Projectors}

In an ordinary three-dimensional vector space, a projector removes every component perpendicular to a chosen subspace. A one-dimensional projector drops a vector orthogonally onto a line; a two-dimensional projector drops it onto a plane. If \(P\) projects onto a specified subspace,
\begin{equation}
  P\ket{\alpha}=\ket{\alpha_P}.
\end{equation}

\begin{figure}[htbp]
\centering
\includegraphics[width=0.76\linewidth]{lecture_05_figure_04.png}
\caption{A geometric projection and its operator notation are developed side by side.}
\lectureframe{00:49:53}
\end{figure}

\begin{figure}[htbp]
\centering
\begin{tikzpicture}[scale=1.05,>=Latex]
  \draw[->] (-0.25,0)--(4.0,0);
  \draw[thick] (-0.25,-0.25)--(3.55,2.6);
  \draw[->,very thick] (0,0)--(2.25,3.05)
    node[above] {\(\ket{\alpha}\)};
  \coordinate (P) at (2.72,1.98);
  \draw[dashed] (2.25,3.05)--(P);
  \draw[->,very thick] (0,0)--(P)
    node[right] {\(\ket{\alpha_P}\)};
  \node[below right] at (3.1,2.25) {projected subspace};
\end{tikzpicture}
\caption{A clean reconstruction of the projection geometry.}
\end{figure}

In the one-spin Hilbert space, the proposition that \(\sigma_3=+1\) corresponds to the one-dimensional subspace spanned by \(\ket{u}\). Its projector acts as
\begin{equation}
  P_{\sigma_3=+1}
  \begin{pmatrix}\alpha\\\beta\end{pmatrix}
  =
  \begin{pmatrix}\alpha\\0\end{pmatrix},
  \qquad
  P_{\sigma_3=-1}
  \begin{pmatrix}\alpha\\\beta\end{pmatrix}
  =
  \begin{pmatrix}0\\\beta\end{pmatrix}.
  \label{eq:l05-projector-actions}
\end{equation}

\begin{classroomqa}
\lecturetimestamp{00:49:55}
\audiencequestion{Does the first projector report how high the vector reaches on a vertical axis?}
\lecturerresponse{No. It simply discards the orthogonal coefficient \(\beta\). The components are Hilbert-space amplitudes, not ordinary spatial coordinates of the spin.}
\end{classroomqa}

\begin{figure}[htbp]
\centering
\includegraphics[width=0.78\linewidth]{lecture_05_figure_05.png}
\caption{Projectors first act on a general spinor and are then written as explicit matrices.}
\lectureframe{00:51:40}
\end{figure}

The matrix form is
\begin{equation}
  P_{\sigma_3=+1}
  =
  \begin{pmatrix}1&0\\0&0\end{pmatrix}
  =
  \frac{I+\sigma_3}{2}.
  \label{eq:l05-p3}
\end{equation}
The same construction works along any unit vector \(\hat n\):
\begin{equation}
  P_{\boldsymbol{\sigma}\cdot\hat n=+1}
  =
  \frac{I+\boldsymbol{\sigma}\cdot\hat n}{2}.
  \label{eq:l05-pn}
\end{equation}

\begin{classroomqa}
\lecturetimestamp{00:53:09}
\audiencequestion{What is the projector onto \(\sigma_1=+1\)?}
\lecturerresponse{Replace \(\sigma_3\) by \(\sigma_1\) in Equation~\eqref{eq:l05-p3}: \(P_{\sigma_1=+1}=(I+\sigma_1)/2\).}
\end{classroomqa}

The Born rule now has a useful proposition form:
\begin{equation}
  P(\mathcal A\text{ is true in }\ket{\psi})
  =
  \bra{\psi}P_{\mathcal A}\ket{\psi}.
  \label{eq:l05-projector-born}
\end{equation}
For a normalized spinor \(\ket{\psi}=(\alpha,\beta)^T\), Equation~\eqref{eq:l05-projector-born} with \(P_{\sigma_3=+1}\) gives \(|\alpha|^2\), exactly as it should.

This language exposes the difference between classical and quantum propositions. In classical logic, a property corresponds to a subset of a set. In quantum mechanics, a property corresponds to a subspace of Hilbert space, or equivalently to its projector. Subsets and vector subspaces do not obey the same lattice operations; Bell's contradiction grows from that difference.

\subsection{The Joint Bell Projector}

For \(A\wedge\neg B\), electron 1 must be up along the third axis and electron 2 must be up along
\begin{equation}
  \hat n_{45}=\frac{\hat x+\hat z}{\sqrt2}.
\end{equation}
The relevant component of the second spin is
\begin{equation}
  \boldsymbol{\tau}\cdot\hat n_{45}
  =
  \frac{\tau_1+\tau_3}{\sqrt2}.
\end{equation}
The joint projector is therefore
\begin{equation}
  P_{A\wedge\neg B}
  =
  \frac{I+\sigma_3}{2}
  \frac{I+(\tau_1+\tau_3)/\sqrt2}{2}.
  \label{eq:l05-joint-projector}
\end{equation}

\begin{classroomqa}
\lecturetimestamp{00:59:18}
\audiencequestion{Why is there a factor \(1/\sqrt2\) in the \(45^\circ\) spin component?}
\lecturerresponse{The direction must be a unit vector. Its first- and third-axis components are both \(1/\sqrt2\), so the corresponding observable is \((\tau_1+\tau_3)/\sqrt2\).}
\end{classroomqa}

\begin{classroomqa}
\lecturetimestamp{01:00:52}
\audiencequestion{Is the entire expression \(I+\boldsymbol{\tau}\cdot\hat n_{45}\) divided by two?}
\lecturerresponse{Yes. The projector is \(\bigl(I+\boldsymbol{\tau}\cdot\hat n_{45}\bigr)/2\); the factor \(1/2\) multiplies the whole numerator.}
\end{classroomqa}

The two factors in Equation~\eqref{eq:l05-joint-projector} act on different tensor factors and commute. Applying both first removes every component with electron 1 down and then every surviving component with electron 2 down along \(\hat n_{45}\).

\section{The Quantum Numbers Reverse the Inequality}

Let
\begin{equation}
  \ket{\psi_s}
  =
  \frac{\ket{ud}-\ket{du}}{\sqrt2}
\end{equation}
denote the singlet. Begin with
\begin{equation}
  P(A,\neg B)
  =
  \bra{\psi_s}P_{A\wedge\neg B}\ket{\psi_s}.
\end{equation}
The factor \((I+\sigma_3)/2\) kills \(\ket{du}\) and leaves \(\ket{ud}/\sqrt2\). The second factor then gives
\begin{align}
  \frac{I+(\tau_1+\tau_3)/\sqrt2}{2}\ket{ud}
  &\notag\\[-0.3em]
  &=
  \frac12
  \left[
    \ket{ud}
    +\frac{\ket{uu}-\ket{ud}}{\sqrt2}
  \right]\nonumber\\
  &=
  \frac12
  \Biggl[
    \left(1-\frac1{\sqrt2}\right)\ket{ud}
  \nonumber\\
  &\qquad
    +\frac1{\sqrt2}\ket{uu}\Biggr].
  \label{eq:l05-projector-action}
\end{align}
The \(\ket{uu}\) term has zero overlap with the singlet bra, and the surviving normalization factors yield
\begin{equation}
  P(A,\neg B)
  =
  \frac14\left(1-\frac1{\sqrt2}\right).
  \label{eq:l05-pab}
\end{equation}

\begin{figure}[htbp]
\centering
\includegraphics[width=0.88\linewidth]{lecture_05_figure_09.png}
\caption{The projector calculation keeps the \(\ket{ud}\) overlap and discards the orthogonal \(\ket{uu}\) contribution.}
\lectureframe{01:08:20}
\end{figure}

Rotational invariance gives the same result for the next pair of directions:
\begin{equation}
  P(B,\neg C)
  =
  \frac14\left(1-\frac1{\sqrt2}\right).
\end{equation}
For \(A\wedge\neg C\), the second axis is the first spatial axis, so
\begin{equation}
  P_{A\wedge\neg C}
  =
  \frac{I+\sigma_3}{2}\frac{I+\tau_1}{2}.
\end{equation}
The first factor again leaves \(\ket{ud}/\sqrt2\). The \(\tau_1\) term produces \(\ket{uu}\), which is orthogonal to the singlet bra, while the identity term survives. Hence
\begin{equation}
  P(A,\neg C)=\frac14.
  \label{eq:l05-pac}
\end{equation}
Combining the three results,
\begin{equation}
  \underbrace{
  P(A,\neg B)+P(B,\neg C)
  }_{\frac12(1-1/\sqrt2)\approx0.146}
  <
  \underbrace{
  P(A,\neg C)
  }_{\frac14=0.25}.
  \label{eq:l05-bell-violation}
\end{equation}

\begin{figure}[htbp]
\centering
\includegraphics[width=0.84\linewidth]{lecture_05_figure_10.png}
\caption{The completed board calculation displays the smaller Bell left-hand side and the larger right-hand side.}
\lectureframe{01:12:15}
\end{figure}

Equation~\eqref{eq:l05-bell-violation} reverses Equation~\eqref{eq:l05-bell-probability}. The conclusion is not that quantum mechanics contradicts itself. It is that the singlet correlations cannot arise from a local classical model in which all three spin propositions possess ordinary set-theoretic truth values independent of the distant measurement choice. A classical simulation can reproduce the correlations only by abandoning at least one Bell assumption, most conspicuously locality.

\section{What the Bell Experiment Tests}

Quantum mechanics already predicts the violation. An experiment must additionally close the classical escape route in which one detector informs the other what setting or outcome it encountered. The two wings should be far enough apart, and the choices and detections close enough in time, that a light-speed signal cannot travel between them. Aspect's experiment made that timing and isolation operationally convincing.

\begin{classroomqa}
\lecturetimestamp{01:13:15}
\audiencequestion{If quantum mechanics already gives the answer, did the Bell experiments teach us anything?}
\lecturerresponse{They tested whether nature follows the quantum prediction in precisely this correlation experiment and, crucially, made the two measurement wings independent on the relevant light-travel time. The experimental control and timing were themselves substantial achievements.}
\end{classroomqa}

\begin{classroomqa}
\lecturetimestamp{01:18:19}
\audiencequestion{What did Bell add after Einstein's discussion of spooky action?}
\lecturerresponse{Bell supplied a quantitative inequality. He identified a feasible correlation experiment whose quantum prediction rules out a local classical underlying account, rather than leaving the objection at the level of philosophical unease. Bell remained ambivalent about whether this was inevitable or profound and became increasingly willing to question standard quantum mechanics.}
\end{classroomqa}

The particles may be miles apart while still sharing one joint state, but this remains a quantum experiment on small systems. Environmental interactions tend to destroy the required phase relations. Air molecules, stray photons, and thermal disturbances can entangle with the pair, so separation must be carried out in a cold, well-isolated environment.

\begin{classroomqa}
\lecturetimestamp{01:19:53}
\audiencequestion{Does a common state spread over a large distance make this a macroscopic quantum object?}
\lecturerresponse{It is better described as small quantum systems far apart. Their separation can be large, but preserving their joint wave function becomes increasingly delicate because uncontrolled interaction with the environment destroys the useful coherence. Two ordinary bowling balls do not provide the same experiment.}
\end{classroomqa}

\begin{classroomqa}
\lecturetimestamp{01:21:23}
\audiencequestion{Can measuring one particle reveal complementary information about the other and thereby evade the uncertainty principle?}
\lecturerresponse{One may measure, for example, the position of one particle and the momentum of the other. That is not a simultaneous measurement of noncommuting observables on one particle. The uncertainty restriction on a single system has not been evaded.}
\end{classroomqa}

\section{When Does a Product Mean And?}

The Bell calculation multiplied two projectors. This represents a conjunction only for compatible propositions.

\begin{classroomqa}
\lecturetimestamp{01:22:39}
\audiencequestion{Why did the product of two projectors represent the logical word and?}
\lecturerresponse{Because the two projectors commute. One acts on electron 1 and the other on electron 2, so their order is irrelevant and both properties can be tested together. For a single electron, \((I+\sigma_3)/2\) and \((I+\sigma_1)/2\) do not commute; spin up along both axes is not a jointly definite proposition of this kind.}
\end{classroomqa}

This is another place where quantum propositions depart from ordinary subsets. A classical intersection always exists. The quantum analogue represented by a product projector has the intended logical meaning only when the relevant subspaces are compatible.

The same discussion points toward measurement itself. A measuring device has a pointer state. If an up electron drives the pointer one way and a down electron drives it another, then a superposed electron becomes entangled with two pointer positions. Measurement may therefore be studied as entanglement between system and apparatus. That problem is postponed in favor of a shorter theorem that uses only linearity.

\begin{classroomqa}
\lecturetimestamp{01:25:09}
\audiencequestion{Is wave-function collapse usefully viewed as one measurement constraining the joint state?}
\lecturerresponse{That question belongs to the measurement process. The apparatus and measured system become entangled: one spin value correlates with one pointer position and the other value with another. The detailed measurement discussion is deferred.}
\end{classroomqa}

\section{Linearity Forbids a Universal Xerox Machine}

Only one fact about time evolution is needed. If a physical process sends
\begin{equation}
  \ket{a}\longrightarrow\ket{a'},
  \qquad
  \ket{b}\longrightarrow\ket{b'},
\end{equation}
then linearity requires
\begin{equation}
  \alpha\ket{a}+\beta\ket{b}
  \longrightarrow
  \alpha\ket{a'}+\beta\ket{b'}.
  \label{eq:l05-linearity}
\end{equation}

An ordinary copying machine can reproduce classical information with controllably small errors. The quantum demand is stronger: one fixed device must duplicate the complete amplitudes of any state presented to it. This is why no-cloning is also called the no-quantum-Xerox principle.

Imagine a machine with one input port and two output ports. If it clones the third-axis basis, it must perform
\begin{equation}
  \ket{u}\longrightarrow\ket{uu},
  \qquad
  \ket{d}\longrightarrow\ket{dd}.
  \label{eq:l05-clone-basis}
\end{equation}
The machine can supply energy, mass, or charge to its outputs; conservation bookkeeping is not the obstruction. The contradiction concerns the state dependence of the transformation.

\begin{figure}[H]
\centering
\includegraphics[width=0.86\linewidth]{lecture_05_figure_11.png}
\caption{A hypothetical one-input, two-output machine is required to duplicate both basis states.}
\lectureframe{01:33:45}
\end{figure}

Now input a spin pointing right,
\begin{equation}
  \ket{r}
  =
  \frac{\ket{u}+\ket{d}}{\sqrt2}.
  \label{eq:l05-right-state}
\end{equation}
Equations~\eqref{eq:l05-linearity} and \eqref{eq:l05-clone-basis} force the output
\begin{equation}
  \ket{r}
  \longrightarrow
  \frac{\ket{uu}+\ket{dd}}{\sqrt2}.
  \label{eq:l05-linear-output}
\end{equation}
Exact cloning instead requires two independent copies:
\begin{align}
  \ket{rr}
  &=
  \ket{r}\otimes\ket{r}\nonumber\\
  &=
  \frac12
  \left(
    \ket{uu}+\ket{ud}+\ket{du}+\ket{dd}
  \right).
  \label{eq:l05-cloned-output}
\end{align}
The mixed terms in Equation~\eqref{eq:l05-cloned-output} are absent from Equation~\eqref{eq:l05-linear-output}. The two outputs are different.

\begin{figure}[H]
\centering
\includegraphics[width=0.72\linewidth]{lecture_05_figure_12.png}
\caption{Linearity produces only the parallel-spin terms, whereas two true right-pointing copies also contain both mixed terms.}
\lectureframe{01:35:20}
\end{figure}

\begin{theorem}[No-cloning]
No single physical machine can exactly clone every unknown quantum state.
\end{theorem}

\begin{proof}
Assume one machine clones the orthogonal basis states as in Equation~\eqref{eq:l05-clone-basis}. Linearity then fixes its output on their superposition to be Equation~\eqref{eq:l05-linear-output}. Exact cloning of the same input demands Equation~\eqref{eq:l05-cloned-output}. Since the two state vectors differ, the assumed universal machine cannot exist.
\end{proof}

The reason is structural. Sending \(\ket{\psi}\) to \(\ket{\psi}\otimes\ket{\psi}\) is quadratic in the amplitudes, whereas every allowed evolution of the complete quantum state is linear.

\begin{classroomqa}
\lecturetimestamp{01:35:49}
\audiencequestion{Does the proof assume that the two outputs are independent?}
\lecturerresponse{Yes, because two clones must be independently usable copies. Independent systems combine by the tensor product, which is why two right-pointing copies are \(\ket{r}\otimes\ket{r}\). Correlated outputs that cannot be used separately are not two clones.}
\end{classroomqa}

\begin{classroomqa}
\lecturetimestamp{01:37:20}
\audiencequestion{A laser seems to copy a photon's state through stimulated emission. Where does the argument fail?}
\lecturerresponse{The atom and the electromagnetic modes also participate. Their state contributes to the output and appears as noise from the standpoint of universal cloning. Approximate cloning is possible, but perfect state-independent cloning is not.}
\end{classroomqa}

\begin{classroomqa}
\lecturetimestamp{01:40:24}
\audiencequestion{Can the quantum information be transferred even though it cannot be duplicated?}
\lecturerresponse{Yes. Transfer is a different operation: the original need not retain the state. No-cloning forbids ending with two independent perfect copies, not moving the state from one carrier to another.}
\end{classroomqa}

\begin{classroomqa}
\lecturetimestamp{01:41:03}
\audiencequestion{Must the input state be unknown?}
\lecturerresponse{If a state is known in advance, one can prepare more systems in that state, or build a device adapted to a chosen orthogonal basis. What is impossible is one machine that accepts an arbitrary input and duplicates it without prior knowledge of which state arrived.}
\end{classroomqa}

\begin{classroomqa}
\lecturetimestamp{01:42:35}
\audiencequestion{Which of the two right-state expressions is the actual machine output?}
\lecturerresponse{Linearity forces Equation~\eqref{eq:l05-linear-output}; cloning demands Equation~\eqref{eq:l05-cloned-output}. Their disagreement is precisely the contradiction.}
\end{classroomqa}

\begin{classroomqa}
\lecturetimestamp{01:43:40}
\audiencequestion{Can the extra information carried by entangled two-particle states be quantified?}
\lecturerresponse{A pure state of one spin has two real physical parameters. Two separately prepared spins supply four, whereas a general pure state in the four-dimensional two-spin Hilbert space has six after normalization and overall phase are removed. The additional states include entangled states, and their degree of entanglement can be quantified by entanglement entropy.}
\end{classroomqa}

\enlargethispage{4\baselineskip}
The sign in Equation~\eqref{eq:l05-two-signs} has now done three jobs. It selected the unique total-spin-zero state, supplied correlations that violate a classical set inequality, and exposed the difference between a linear quantum evolution and the nonlinear act of making two copies. Each result comes from treating the joint state as a vector in its full tensor-product space rather than as a list of pre-existing local properties.
